\chapter{2011年新课标卷（文）}

\section{单选题}

\begin{problem}
已知集合 \(M = \{0, 1, 2, 3, 4\}\)，\(N = \{1, 3, 5\}\)，\(P = M \cap N\)，则 \(P\) 的子集共有（\quad）

\choices
    {\(2\) 个}
    {\(4\) 个}
    {\(6\) 个}
    {\(8\) 个}

\end{problem}

\begin{answer}
B
\end{answer}

\begin{solution}
由 \(M\cap N=\{1,3\}\)，集合 \(P\) 有 \(2\) 个元素，所以其子集共有 \(2^2=4\) 个，故选 B.
\end{solution}

\begin{problem}
复数 \(\frac{5\i}{1 - 2\i} =\)（\quad）

\choices
    {\(2 - \i\)}
    {\(1 - 2\i\)}
    {\(-2 + \i\)}
    {\(-1 + 2\i\)}

\end{problem}

\begin{answer}
C
\end{answer}

\begin{solution}
分子、分母同乘 \(1+2\i\)，得 \(\frac{5\i}{1-2\i}=\frac{5\i(1+2\i)}{5}=\i(1+2\i)=-2+\i\)，故选 C.
\end{solution}

\begin{problem}
下列函数中，既是偶函数又在 \((0, +\infty)\) 单调递增的函数是（\quad）

\choices
    {\(y = x^{3}\)}
    {\(y = |x| + 1\)}
    {\(y = -x^{2} + 1\)}
    {\(y = 2^{-|x|}\)}

\end{problem}

\begin{answer}
B
\end{answer}

\begin{solution}
函数 \(y=x^3\) 是奇函数，函数 \(y=-x^2+1\) 和 \(y=2^{-|x|}\) 在 \((0,+\infty)\) 上均单调递减. 函数 \(y=|x|+1\) 是偶函数，且在 \((0,+\infty)\) 上单调递增，故选 B.
\end{solution}

\begin{problem}
椭圆 \(\frac{x^2}{16} + \frac{y^2}{8} = 1\) 的离心率为（\quad）

\choices
    {\( \frac{1}{3} \)}
    {\( \frac{1}{2} \)}
    {\( \frac{\sqrt{3}}{3} \)}
    {\(\frac{\sqrt{2}}{2}\)}

\end{problem}

\begin{answer}
D
\end{answer}

\begin{solution}
椭圆的长半轴、短半轴分别为 \(a=4\)，\(b=2\sqrt{2}\)，所以 \(c=\sqrt{a^2-b^2}=\sqrt{8}=2\sqrt{2}\). 因此离心率为 \(e=\frac{c}{a}=\frac{\sqrt{2}}{2}\)，故选 D.
\end{solution}

\begin{problem}
执行如图的程序框图，如果输入的 \(N\) 是 6，那么输出的 \(p\) 是（\quad）

\bitmapfigure[max width=0.48\linewidth]{img/2011/new_curriculum_liberal/q5_flowchart.png}

\choices
    {\(120\)}
    {\(720\)}
    {\(1440\)}
    {\(5040\)}

\end{problem}

\begin{answer}
A
\end{answer}

\begin{solution}
初始时 \(k=1\)，\(p=1\). 每次先令 \(p\leftarrow pk\)，再令 \(k\leftarrow k+1\)，当 \(k<N\) 时继续循环. 输入 \(N=6\) 时，循环中的乘法依次使 \(p\) 变为 \(1\)，\(2\)，\(6\)，\(24\)，\(120\)，此时 \(k=6\)，不满足 \(k<6\)，所以输出 \(p=120\)，故选 A.
\end{solution}

\begin{problem}
有 \(3\) 个兴趣小组，甲，乙两位同学各自参加其中一个小组，每位同学参加各个小组的可能性相同，则这两位同学参加同一个兴趣小组的概率为（\quad）

\choices
    {\( \frac{1}{3} \)}
    {\( \frac{1}{2} \)}
    {\( \frac{2}{3} \)}
    {\( \frac{3}{4} \)}

\end{problem}

\begin{answer}
A
\end{answer}

\begin{solution}
甲、乙两位同学各有 \(3\) 种选择，共有 \(3\times3=9\) 种等可能结果. 两人参加同一小组有 \(3\) 种结果，因此所求概率为 \(\frac{3}{9}=\frac{1}{3}\)，故选 A.
\end{solution}

\begin{problem}
已知角 \(\theta\) 的顶点与原点重合，始边与 \(x\) 轴的正半轴重合，终边在直线 \(y = 2x\) 上，则 \(\cos 2\theta =\)（\quad）

\choices
    {\(-\frac{1}{5}\)}
    {\(-\frac{3}{5}\)}
    {\( \frac{3}{5} \)}
    {\( \frac{4}{5} \)}

\end{problem}

\begin{answer}
B
\end{answer}

\begin{solution}
终边在直线 \(y=2x\) 上时，取终边上的点 \((1,2)\)，则 \(\tan\theta=2\). 由二倍角公式，
\(\cos2\theta=\frac{1-\tan^2\theta}{1+\tan^2\theta}=\frac{1-4}{1+4}=-\frac{3}{5}\)，故选 B.
\end{solution}

\begin{problem}
在一个几何体的三视图中，正视图和俯视图如下图所示，则相应的侧视图可以为（\quad）

\bitmapfigure[max width=0.42\linewidth]{img/2011/new_curriculum_liberal/q8_condition.png}

\choices
    {\choicebitmap[width=0.15\paperwidth]{img/2011/new_curriculum_liberal/q8_optA.png}}
    {\choicebitmap[width=0.15\paperwidth]{img/2011/new_curriculum_liberal/q8_optB.png}}
    {\choicebitmap[width=0.15\paperwidth]{img/2011/new_curriculum_liberal/q8_optC.png}}
    {\choicebitmap[width=0.15\paperwidth]{img/2011/new_curriculum_liberal/q8_optD.png}}

\end{problem}

\begin{answer}
D
\end{answer}

\begin{solution}
正视图是以中线为对称轴的三角形，俯视图由上方的半圆和下方的三角形组成，因此几何体沿侧视方向的外轮廓仍应是等腰三角形. 俯视图中直径中点对应的竖直棱在侧视图中是可见棱，应画成实线. A、B 的外轮廓是梯形，C 虽为三角形但中线为虚线，只有 D 同时满足外轮廓和棱的可见性，故选 D.
\end{solution}

\begin{problem}
已知直线 \(l\) 过抛物线 \(C\) 的焦点，且与 \(C\) 的对称轴垂直，\(l\) 与 \(C\) 交于 \(A\)，\(B\) 两点，\( \left|AB\right|=12 \)，\(P\) 为 \(C\) 的准线上一点，则 \( \triangle ABP \) 的面积为（\quad）

\choices
    {\(18\)}
    {\(24\)}
    {\(36\)}
    {\(48\)}

\end{problem}

\begin{answer}
C
\end{answer}

\begin{solution}
设抛物线方程为 \(y^2=2px\)（\(p>0\)），则其焦点为 \(\left(\frac{p}{2},0\right)\)，准线为 \(x=-\frac{p}{2}\). 过焦点且垂直于对称轴的直线为通径所在直线，故 \(\left|AB\right|=2p=12\)，从而 \(p=6\). 点 \(P\) 在准线上，故点 \(P\) 到直线 \(AB\) 的距离为 \(p=6\)，所以 \(S_{\triangle ABP}=\frac{1}{2}\times12\times6=36\)，故选 C.
\end{solution}

\begin{problem}
在下列区间中，函数 \( f(x) = \e^{x} + 4x - 3 \) 的零点所在的区间为（\quad）

\choices
    {\(\left(-\frac{1}{4}, 0\right)\)}
    {\(\left(0, \frac{1}{4}\right)\)}
    {\(\left(\frac{1}{4}, \frac{1}{2}\right)\)}
    {\(\left(\frac{1}{2}, \frac{3}{4}\right)\)}
\end{problem}

\begin{answer}
C
\end{answer}

\begin{solution}
函数 \(f(x)=\e^x+4x-3\) 的导数为 \(f'(x)=\e^x+4>0\)，所以 \(f(x)\) 在 \(\R\) 上单调递增. 又 \(1<\e<16\)，故
\(f\left(\frac{1}{4}\right)=\e^{1/4}-2<0\)，而
\(f\left(\frac{1}{2}\right)=\e^{1/2}-1>0\). 由零点存在性定理，零点位于 \(\left(\frac{1}{4},\frac{1}{2}\right)\) 内；由单调性知零点唯一，故选 C.
\end{solution}

\begin{problem}
设函数 \( f(x) = \sin \left(2x + \frac{\pi}{4}\right) + \cos \left(2x + \frac{\pi}{4}\right) \)，则（\quad）

\choices
    {\(y = f(x)\) 在 \(\left(0, \frac{\pi}{2}\right)\) 单调递增，其图象关于直线 \(x = \frac{\pi}{4}\) 对称}
    {\(y = f(x)\) 在 \(\left(0, \frac{\pi}{2}\right)\) 单调递增，其图象关于直线 \(x = \frac{\pi}{2}\) 对称}
    {\(y = f(x)\) 在 \(\left(0, \frac{\pi}{2}\right)\) 单调递减，其图象关于直线 \(x = \frac{\pi}{4}\) 对称}
    {\( y = f(x) \) 在 \( \left(0, \frac{\pi}{2}\right) \) 单调递减，其图象关于直线 \( x = \frac{\pi}{2} \) 对称}

\end{problem}

\begin{answer}
D
\end{answer}

\begin{solution}
由和角公式，\(f(x)=\sqrt{2}\sin\left(2x+\frac{\pi}{2}\right)=\sqrt{2}\cos2x\). 当 \(0<x<\frac{\pi}{2}\) 时，\(0<2x<\pi\)，故 \(f'(x)=-2\sqrt{2}\sin2x<0\)，函数单调递减；又 \(f\left(\frac{\pi}{2}+t\right)=\sqrt{2}\cos(\pi+2t)=f\left(\frac{\pi}{2}-t\right)\)，图象关于直线 \(x=\frac{\pi}{2}\) 对称，故选 D.
\end{solution}

\begin{problem}
已知函数 \( y = f(x) \) 的周期为 2，当 \( x \in [-1, 1] \) 时 \( f(x) = x^2 \)，那么函数 \( y = f(x) \) 的图象与函数 \( y = | \lg x | \) 的图象的交点共有（\quad）

\choices
    {\(10\) 个}
    {\(9\) 个}
    {\(8\) 个}
    {\(1\) 个}

\end{problem}

\begin{answer}
A
\end{answer}

\begin{solution}
对 \(x>0\)，由周期性知，在区间 \([2k-1,2k+1]\) 上有 \(f(x)=(x-2k)^2\). 先看 \(0<x<1\)，此时 \(f(x)=x^2\)，方程化为 \(x^2+\lg x=0\). 左端函数导数为 \(2x+\frac{1}{x\ln10}>0\)，且当 \(x\to0^+\) 时趋于负无穷，在 \(x=1\) 时为 \(1\)，故有且只有一个交点.

对 \(k=1\)，\(2\)，\(3\)，\(4\)，在 \([2k-1,2k]\) 上，\((x-2k)^2\) 递减而 \(\lg x\) 递增，端点处两者之差分别为 \(1-\lg(2k-1)>0\) 和 \(-\lg(2k)<0\)，故各有一个交点. 在 \([2k,2k+1]\) 上，令 \(g_k(x)=(x-2k)^2-\lg x\)，则 \(g_k''(x)=2+\frac{1}{x^2\ln10}>0\)，且 \(g_k(2k)<0\)、\(g_k(2k+1)=1-\lg(2k+1)>0\)，故各有一个交点.

当 \(k=5\) 时，左半段 \([9,10]\) 有一个交点；右半段 \([10,11]\) 两端的差分别为 \(-1\) 和 \(1-\lg11<0\)，结合 \(g_5''(x)>0\) 可知该段无交点. 若 \(x\geqslant11\)，则 \(|\lg x|=\lg x>1\)，而 \(0\leqslant f(x)\leqslant1\)，也无交点. 因此交点总数为 \(1+2\times4+1=10\)，故选 A.
\end{solution}

\section{填空题}

\begin{problem}
已知 \(\bs{a}\) 与 \(\bs{b}\) 为两个不共线的单位向量，\(k\) 为实数，若向量 \(\bs{a} + \bs{b}\) 与向量 \(k\bs{a} - \bs{b}\) 垂直，则 \(k = \)\fillinblank{}.

\end{problem}

\begin{answer}
1
\end{answer}

\begin{solution}
由垂直条件，\((\bs{a}+\bs{b})\cdot(k\bs{a}-\bs{b})=0\). 设 \(\bs{a}\cdot\bs{b}=c\). 因为 \(\bs{a}\)，\(\bs{b}\) 是不共线的单位向量，所以 \(c=\cos\langle\bs{a},\bs{b}\rangle\)，且 \(-1<c<1\). 因此
\[
(\bs{a}+\bs{b})\cdot(k\bs{a}-\bs{b})=k-c+kc-1=(k-1)(1+c)
\]
由 \(1+c>0\)，垂直条件化为 \(k-1=0\)，故 \(k=1\).
\end{solution}

\begin{problem}
若变量 \(x\)，\(y\) 满足约束条件 \(\begin{cases}
3 \leqslant 2x + y \leqslant 9,\\
6 \leqslant x - y \leqslant 9,
\end{cases}\) 则 \(z = x + 2y\) 的最小值为\fillinblank{}.
\end{problem}

\begin{answer}
\(-6\)
\end{answer}

\begin{solution}
令 \(u=2x+y\)，\(v=x-y\)，则 \(3\leqslant u\leqslant9\)，\(6\leqslant v\leqslant9\). 由 \(x=\frac{u+v}{3}\)，\(y=\frac{u-2v}{3}\)，得 \(z=x+2y=u-v\). 因此 \(z\) 在 \(u=3\)，\(v=9\) 时取得最小值 \(3-9=-6\).
\end{solution}

\begin{problem}
\(\triangle ABC\) 中，\(B = 120^{\circ}\)，\(AC = 7\)，\(AB = 5\)，则 \(\triangle ABC\) 的面积为\fillinblank{}.

\end{problem}

\begin{answer}
\(\frac{15\sqrt{3}}{4}\)
\end{answer}

\begin{solution}
设 \(BC=a\). 由余弦定理，
\[
AC^2=AB^2+BC^2-2\cdot AB\cdot BC\cos B
  =25+a^2+5a
\]
所以 \(a^2+5a-24=0\)，即 \((a-3)(a+8)=0\). 由于 \(a>0\)，得 \(a=3\). 于是
\[
S_{\triangle ABC}
 =\frac{1}{2}\cdot5\cdot3\cdot\frac{\sqrt{3}}{2}
 =\frac{15\sqrt{3}}{4}
\]
\end{solution}

\begin{problem}
已知两个圆锥有公共底面，且两圆锥的顶点和底面的圆周都在同一个球面上. 若圆锥底面面积是这个球面面积的 \(\frac{3}{16}\)，则这两个圆锥中，体积较小者的高与体积较大的高的比值为\fillinblank{}.

\end{problem}

\begin{answer}
\(\frac{1}{3}\)
\end{answer}

\begin{solution}
设球心为 \(O\)，公共底面圆心为 \(H\)，球的半径为 \(R\)，公共底面半径为 \(r\). 由于底面圆周在球面上，\(OH\) 垂直于底面，且

\[
r^2+OH^2=R^2
\]

由题意，\(\pi r^2=\frac{3}{16}\cdot4\pi R^2=\frac{3}{4}\pi R^2\)，所以 \(r=\frac{\sqrt{3}}{2}R\). 设 \(d=OH\)，则
\[
d=\sqrt{R^2-r^2}=\frac{R}{2}
\]
两个圆锥的顶点分别是底面轴线与球面在两侧的交点，所以它们的高分别为 \(R-d=\frac{R}{2}\) 和 \(R+d=\frac{3R}{2}\). 圆锥底面相同，体积大小规律与高大小规律相同，故所求比值为
\[
\frac{R/2}{3R/2}=\frac{1}{3}
\]
\end{solution}

\section{解答题}

\begin{problem}
已知等比数列 \(\{a_{n}\}\) 中，\(a_{1} = \frac{1}{3}\)，公比 \(q = \frac{1}{3}\).
\begin{enumerate}
\item 若 \(S_{n}\) 为 \(\{a_{n}\}\) 的前 \(n\) 项和，证明： \(S_{n} = \frac{1 - a_{n}}{2}\)；
\item 设 \(b_{n} = \log_{3} a_{1} + \log_{3} a_{2} + \cdots + \log_{3} a_{n}\)，求数列 \(\{b_{n}\}\) 的通项公式.
\end{enumerate}
\end{problem}

\begin{solution}
\begin{enumerate}
\item 由等比数列通项公式，\(a_n=a_1q^{n-1}=\frac{1}{3^n}\). 因此
\[
S_n=\frac{a_1(1-q^n)}{1-q}
 =\frac{\frac{1}{3}\left(1-\frac{1}{3^n}\right)}{1-\frac{1}{3}}
 =\frac{1-\frac{1}{3^n}}{2}
 =\frac{1-a_n}{2}
\]
\item 由 \(a_i=3^{-i}\)，得 \(\log_3a_i=-i\). 所以
\[
b_n=\sum_{i=1}^n\log_3a_i=-\sum_{i=1}^n i=-\frac{n(n+1)}{2}
\]
\end{enumerate}
\end{solution}

\begin{problem}
如图，四棱锥 \(P - ABCD\) 中，底面 \(ABCD\) 为平行四边形，\(\angle DAB = 60^{\circ}\)，\(AB = 2AD\)，\(PD \perp\) 底面 \(ABCD\).
\begin{enumerate}
\item 证明： \(PA \perp BD\)；
\item 设 \(PD = AD = 1\)，求棱锥 \(D - PBC\) 的高.
\end{enumerate}

\bitmapfigure[max width=0.55\linewidth]{img/2011/new_curriculum_liberal/q18_pyramid.png}
\end{problem}

\begin{solution}
\begin{enumerate}
\item 以 \(AD\) 为长度单位，令 \(AD=1\)，可在底面内取坐标
    取 \(A=(0,0,0)\)，\(B=(2,0,0)\)，\(D=\left(\frac{1}{2},\frac{\sqrt{3}}{2},0\right)\).
由于 \(PD\perp\) 底面，设 \(PD=h\)，则 \(P=\left(\frac{1}{2},\frac{\sqrt{3}}{2},h\right)\). 于是
    \(\overrightarrow{PA}=\left(-\frac{1}{2},-\frac{\sqrt{3}}{2},-h\right)\)，
    \(\overrightarrow{BD}=\left(-\frac{3}{2},\frac{\sqrt{3}}{2},0\right)\)，
从而 \(\overrightarrow{PA}\cdot\overrightarrow{BD}=\frac{3}{4}-\frac{3}{4}=0\)，故 \(PA\perp BD\).
\item 由 \(PD=AD=1\)，取 \(h=1\)，则
    \(C=B+D=\left(\frac{5}{2},\frac{\sqrt{3}}{2},0\right)\)，
    \(P=\left(\frac{1}{2},\frac{\sqrt{3}}{2},1\right)\).
    在平面 \(PBC\) 内取向量
    \(\overrightarrow{PB}=\left(\frac{3}{2},-\frac{\sqrt{3}}{2},-1\right)\)，
    \(\overrightarrow{PC}=\left(2,0,-1\right)\).
它们的一个法向量为 \(\bs n=(\sqrt{3},-1,2\sqrt{3})\)，故平面 \(PBC\) 的方程为
\[
\sqrt{3}x-y+2\sqrt{3}z-2\sqrt{3}=0
\]
将 \(D=\left(\frac{1}{2},\frac{\sqrt{3}}{2},0\right)\) 代入，点 \(D\) 到平面 \(PBC\) 的距离为
\[
h_D=\frac{\left|\frac{\sqrt{3}}{2}-\frac{\sqrt{3}}{2}-2\sqrt{3}\right|}{\sqrt{3+1+12}}
=\frac{2\sqrt{3}}{4}
=\frac{\sqrt{3}}{2}
\]
所以棱锥 \(D-PBC\) 的高为 \(\frac{\sqrt{3}}{2}\).
\end{enumerate}
\end{solution}

\begin{problem}
某种产品的质量以其质量指标值衡量，质量指标值越大表明质量越好，且质量指标值大于或等于 \(102\) 的产品为优质品. 现用两种新配方（分别称为 A 配方和 B 配方）做试验，各生产了 \(100\) 件这种产品，并测量了每件产品的质量指标值，得到了下面试验结果：

\begin{center}
A 配方的频数分布表

\begin{tabular}{c|ccccc}
\hline
指标值分组 & \(\left[90,94\right)\) & \(\left[94,98\right)\) & \(\left[98,102\right)\) & \(\left[102,106\right)\) & \(\left[106,110\right]\) \\
\hline
频数 & \(8\) & \(20\) & \(42\) & \(22\) & \(8\) \\
\hline
\end{tabular}

B 配方的频数分布表

\begin{tabular}{c|ccccc}
\hline
指标值分组 & \(\left[90,94\right)\) & \(\left[94,98\right)\) & \(\left[98,102\right)\) & \(\left[102,106\right)\) & \(\left[106,110\right]\) \\
\hline
频数 & \(4\) & \(12\) & \(42\) & \(32\) & \(10\) \\
\hline
\end{tabular}
\end{center}

\begin{enumerate}
\item 分别估计用 A 配方，B 配方生产的产品优质品率；
\item 已知用 B 配方生产一件产品的利润 \(y\)（单位：元）与其质量指标值 \(t\) 的关系式为 \(y = \begin{cases}
-2,  t < 94,\\
2,  94 \leqslant t < 102,\\
4,  t \geqslant 102.
\end{cases}\)
估计用 B 配方生产的一件产品的利润大于 \(0\) 的概率，并求用 B 配方生产的上述 \(100\) 件产品平均一件的利润.
\end{enumerate}
\end{problem}

\begin{solution}
\begin{enumerate}
\item A 配方生产的产品中，质量指标值不小于 \(102\) 的有 \(22+8=30\) 件，所以优质品率估计为 \(\frac{30}{100}=30\%\). B 配方中相应的产品有 \(32+10=42\) 件，所以优质品率估计为 \(\frac{42}{100}=42\%\).
\item B 配方产品的利润大于 \(0\) 当且仅当 \(t\geqslant94\)，相应的产品数为 \(12+42+32+10=96\)，所以所求概率估计为 \(\frac{96}{100}=0.96\). 这 \(100\) 件产品的总利润为
\[
(-2)\times4+2\times(12+42)+4\times(32+10)=268
\]
故平均一件的利润为 \(\frac{268}{100}=2.68\) 元.
\end{enumerate}
\end{solution}

\begin{problem}
在平面直角坐标系 \(xOy\) 中，曲线 \(y = x^{2} - 6x + 1\) 与坐标轴的交点都在圆 \(C\) 上.
\begin{enumerate}
\item 求圆 \(C\) 的方程；
\item 若圆 \(C\) 与直线 \(x - y + a = 0\) 交于 \(A\)，\(B\) 两点，且 \(OA \perp OB\)，求 \(a\) 的值.
\end{enumerate}
\end{problem}

\begin{solution}
\begin{enumerate}
\item 曲线与 \(y\) 轴的交点为 \((0,1)\)，与 \(x\) 轴的交点横坐标满足
\[
x^2-6x+1=0
\]
设圆 \(C\) 的方程为 \(x^2+y^2+Dx+Ey+F=0\). 圆与 \(x\) 轴的两个交点横坐标的和、积分别为 \(6\)、\(1\)，所以 \(D=-6\)，\(F=1\). 代入点 \((0,1)\)，得 \(1+E+1=0\)，即 \(E=-2\). 因此
\[
C:\ x^2+y^2-6x-2y+1=0
\]
\item 将 \(y=x+a\) 代入圆 \(C\)，得
\[
2x^2+(2a-8)x+(a-1)^2=0
\]
设交点 \(A\)，\(B\) 的横坐标分别为 \(x_1\)，\(x_2\)，则
    由韦达定理，\(x_1+x_2=4-a\)，\(x_1x_2=\frac{(a-1)^2}{2}\).
由 \(OA\perp OB\)，有
\[
0=\overrightarrow{OA}\cdot\overrightarrow{OB}
=x_1x_2+(x_1+a)(x_2+a)
=2x_1x_2+a(x_1+x_2)+a^2
=(a-1)^2+a(4-a)+a^2
=(a+1)^2
\]
所以 \(a=-1\).
\end{enumerate}
\end{solution}

\begin{problem}
已知函数 \( f(x) = a \frac{\ln x}{x + 1} + \frac{b}{x} \)，曲线 \( y = f(x) \) 在点 \((1, f(1))\) 处的切线方程为 \( x + 2y - 3 = 0 \).
\begin{enumerate}
\item 求 \(a\)，\(b\) 的值；
\item 证明：当 \(x > 0\)，且 \(x \neq 1\) 时，\(f(x) > \frac{\ln x}{x - 1}\).
\end{enumerate}
\end{problem}

\begin{solution}
\begin{enumerate}
\item 因为 \(\ln1=0\)，所以 \(f(1)=b\). 切线方程为 \(y=\frac{3-x}{2}\)，故 \(f(1)=1\)，从而 \(b=1\). 又
\[
f'(x)=a\frac{\frac{x+1}{x}-\ln x}{(x+1)^2}-\frac{b}{x^2}
\]
因此 \(f'(1)=\frac{a}{2}-1\). 切线斜率为 \(-\frac{1}{2}\)，所以 \(\frac{a}{2}-1=-\frac{1}{2}\)，得 \(a=1\).
\item 此时 \(f(x)=\frac{\ln x}{x+1}+\frac{1}{x}\)，故
\[
f(x)-\frac{\ln x}{x-1}
=\frac{1}{x}-\frac{2\ln x}{x^2-1}
=\frac{2h(x)}{x^2-1}
\]
其中 \(h(x)=\frac{1}{2}\left(x-\frac{1}{x}\right)-\ln x\). 有
\(h'(x)=\frac{(x-1)^2}{2x^2}\geqslant0\)，\(h(1)=0\).
所以当 \(x>1\) 时 \(h(x)>0\)，当 \(0<x<1\) 时 \(h(x)<0\)；而 \(x^2-1\) 与 \(h(x)\) 同号. 因此当 \(x>0\) 且 \(x\neq1\) 时，\(\frac{2h(x)}{x^2-1}>0\)，即
\[
f(x)>\frac{\ln x}{x-1}
\]
\end{enumerate}
\end{solution}

\begin{problem}
如图，\(D\)，\(E\) 分别为 \(\triangle ABC\) 的边 \(AB\)，\(AC\) 上的点，且不与 \(\triangle ABC\) 的顶点重合.
已知 \(AE\) 的长为 \(m\)，\(AC\) 的长为 \(n\)，\(AD\)，\(AB\) 的长是关于 \(x\) 的方程 \(x^{2} - 14x + mn = 0\) 的两个根.
\begin{enumerate}
\item 证明： \(C\)，\(B\)，\(D\)，\(E\) 四点共圆；
\item 若 \(\angle A = 90^{\circ}\)，且 \(m = 4\)，\(n = 6\)，求 \(C\)，\(B\)，\(D\)，\(E\) 所在圆的半径.
\end{enumerate}

\bitmapfigure[max width=0.48\linewidth]{img/2011/new_curriculum_liberal/q22_triangle.png}
\end{problem}

\begin{solution}
\begin{enumerate}
\item 由 \(AD\)，\(AB\) 是方程 \(x^2-14x+mn=0\) 的两个根，得
\[
AD\cdot AB=mn=AE\cdot AC
\]
由于 \(D\) 在 \(AB\) 上，\(E\) 在 \(AC\) 上，根据相交弦定理的逆定理，四点 \(C\)，\(B\)，\(D\)，\(E\) 共圆.
\item 当 \(\angle A=90^\circ\)，\(m=4\)，\(n=6\) 时，\(AD\)，\(AB\) 是方程
\[
x^2-14x+24=0
\]
的两个根，即 \(AD\)，\(AB\) 为 \(2\)，\(12\). 由 \(D\) 在边 \(AB\) 上，得 \(AD=2\)，\(AB=12\). 取 \(A\) 为原点，令 \(AB\) 在 \(x\) 轴正向、\(AC\) 在 \(y\) 轴正向，则
    \(B=(12,0)\)，\(C=(0,6)\)，\(D=(2,0)\)，\(E=(0,4)\).
设所求圆方程为 \(x^2+y^2+ux+vy+w=0\). 代入 \(B\)，\(D\)，得 \(x\) 轴上的两个根为 \(2,12\)，故 \(u=-14\)，\(w=24\)；代入 \(C\)，\(E\)，得 \(y\) 轴上的两个根为 \(4,6\)，故 \(v=-10\). 所以圆心为 \((7,5)\)，半径满足
\[
R^2=7^2+5^2-24=50
\]
从而 \(R=5\sqrt{2}\).
\end{enumerate}
\end{solution}

\begin{problem}
在直角坐标系 \(xOy\) 中，曲线 \(C_1\) 的参数方程为 \(\begin{cases}
x = 2\cos \alpha,\\
y = 2 + 2\sin \alpha,
\end{cases}\)（\(\alpha\) 为参数），\(M\) 是 \(C_1\) 上的动点，\(P\) 点满足 \(\overrightarrow{OP} = 2\overrightarrow{OM}\)，\(P\) 点的轨迹为曲线 \(C_2\).
\begin{enumerate}
\item 求 \(C_2\) 的方程；
\item 在以 \(O\) 为极点，\(x\) 轴的正半轴为极轴的极坐标系中，射线 \( \theta = \frac{\pi}{3} \) 与 \(C_1\) 的异于极点的交点为 \(A\)，与 \(C_2\) 的异于极点的交点为 \(B\)，求 \(|AB|\).
\end{enumerate}
\end{problem}

\begin{solution}
\begin{enumerate}
\item 由 \(\overrightarrow{OP}=2\overrightarrow{OM}\)，若 \(M=(2\cos\alpha,2+2\sin\alpha)\)，则
\[
P=(4\cos\alpha,4+4\sin\alpha)
\]
消去参数 \(\alpha\)，得
\[
C_2:\ x^2+(y-4)^2=16
\]
\item 射线 \(\theta=\frac{\pi}{3}\) 上的点满足 \(x=r\cos\theta\)，\(y=r\sin\theta\). 代入 \(C_1:x^2+(y-2)^2=4\)，得
\[
r^2-4r\sin\theta=0
\]
除去极点 \(r=0\)，在 \(\theta=\frac{\pi}{3}\) 时 \(r_A=4\sin\frac{\pi}{3}=2\sqrt{3}\). 由于 \(C_2\) 是 \(C_1\) 关于原点的位似放大 \(2\) 倍，故 \(r_B=2r_A=4\sqrt{3}\). 两点在同一射线上，所以
\[
|AB|=r_B-r_A=2\sqrt{3}
\]
\end{enumerate}
\end{solution}

\begin{problem}
设函数 \( f(x) = |x - a| + 3x \)，其中 \( a > 0 \).
\begin{enumerate}
\item 当 \(a = 1\) 时，求不等式 \(f(x) \geqslant 3x + 2\) 的解集；
\item 若不等式 \( f(x) \leqslant 0 \) 的解集为 \( \{x \mid x \leqslant -1\} \)，求 \( a \) 的值.
\end{enumerate}
\end{problem}

\begin{solution}
\begin{enumerate}
\item 当 \(a=1\) 时，不等式化为 \(|x-1|\geqslant2\)，所以 \(x-1\geqslant2\) 或 \(x-1\leqslant-2\)，解集为
\[
(-\infty,-1]\cup[3,+\infty)
\]
\item 当 \(x\geqslant a\) 时，\(f(x)=4x-a\). 条件 \(f(x)\leqslant0\) 要求 \(x\leqslant\frac{a}{4}\)，与 \(x\geqslant a>0\) 不相容. 当 \(x<a\) 时，\(f(x)=a+2x\)，故 \(f(x)\leqslant0\) 等价于 \(x\leqslant-\frac{a}{2}\). 因此不等式的解集为 \(\left(-\infty,-\frac{a}{2}\right]\). 由题设它等于 \(\{x\mid x\leqslant-1\}\)，得 \(-\frac{a}{2}=-1\)，所以 \(a=2\).
\end{enumerate}
\end{solution}
